\SourceLeaf{280}
\begin{list}{}{\leftmargin=2em\rightmargin=0pt\itemindent=0pt\itemsep=0pt\parsep=0pt}
\item erstens die Dreiecke sich nicht allzu weit vom Normalparallelkreis entfernen und
\item zweitens die Dreiecke im Verhältnis ihrer Seiten zu einem ganzen Erdquadranten klein sind, wie es bei tatsächlich messbaren Dreiecken stets der Fall ist.
\end{list}
\noindent\hspace*{\parindent}Die auf die Kugelfläche übertragenen Dreiecksseiten schmiegen sich bei der oben betrachteten konformen Abbildung noch viel enger an Großkreisbögen an als bei der in der genannten Schrift vorgeschlagenen. Bei einem Abstand von \(2\tfrac{1}{2}\) Grad vom Normalparallelkreis ergab jene [nach Art. 13] eine lineare Vergrößerung um \(\tfrac{1}{530000}\). Nach der neuen Methode würde die Änderung nur \(\tfrac{1}{5800000}\) betragen.

Man kann daher zunächst eine Dreiecksseite auf geeignete Weise auf die Kugelfläche übertragen und dann das ganze System mithilfe der Winkel so berechnen, als läge es auf dieser Fläche selbst, nötigenfalls mit der eben angedeuteten Anpassung. Anschließend kann man für alle Punkte die Breiten und Längen bestimmen und von diesen mithilfe der oben angegebenen Formeln, oder genauer, für die Breiten mithilfe einer Hilfstafel wie der hier beigefügten, zu den Breiten und Längen auf der Ellipsoidfläche übergehen.

\Article{12}
\indent Es bleibt also noch zu bestimmen, wie weit eine auf die Kugelfläche übertragene geodätische Linie von dem Großkreisbogen zwischen denselben Endpunkten abweicht. Damit lässt sich zugleich in jedem Fall beurteilen, ob die Abweichung berücksichtigt werden muss. Die Aufgabe lässt sich auf mehr als eine Weise behandeln. Für den vorliegenden Zweck, bei dem die Reduktion immer nur sehr klein ausfallen kann, scheint das folgende Verfahren am besten geeignet.

Es sei \(L\) die betrachtete geodätische Linie auf dem Ellipsoid ohne festgelegte Begrenzung, \(M\) ihre konforme Abbildung auf der Kugelfläche, \(F\) und \(G\) die Endpunkte eines bestimmten Stücks von \(M\) und schließlich \(N\) ein Großkreis durch diese beiden Punkte. Jeder Punkt auf \(N\) werde durch seinen Abstand \(x\) von einem zunächst willkürlich auf \(N\) gewählten Ursprung bestimmt, jeder Punkt von \(M\) durch seinen senkrechten Abstand \(y\) von \(N\) und durch den zum Fußpunkt dieses Lots gehörenden Wert
